\documentclass[12pt]{ximera}
\typeout{Start loading xmPreamble.tex}%

% Add here extra macro's that are loaded automatically by all documents of claas 'ximera' or 'xourse' in this repo

%%
%%  Example:
%%
% \newcommand{\R}{\mathbb{R}}
\usepackage{tikz}
\usetikzlibrary{positioning, arrows.meta}
\usepackage{multicol}
\usepackage{enumitem}
\usepackage{caption}
\usepackage{amsmath}
\usepackage{amsthm}
\usepackage{fontenc}

%% ---------------------------------------------------------------
%% Shared worksheet macros.
%%
%% These came from the Summer 2026 worksheet set, where each worksheet carried its
%% own copy. They have to live here instead: a xourse inlines every activity into a
%% single document, so duplicate \newcommand definitions across activities collide,
%% and the xourse gobbles \input so a per-topic macro file never loads.
%%
%% They use \providecommand, NOT \newcommand. ximera.cls loads this file automatically,
%% and every activity also carries an explicit \typeout{Start loading xmPreamble.tex}%

% Add here extra macro's that are loaded automatically by all documents of claas 'ximera' or 'xourse' in this repo

%%
%%  Example:
%%
% \newcommand{\R}{\mathbb{R}}
\usepackage{tikz}
\usetikzlibrary{positioning, arrows.meta}
\usepackage{multicol}
\usepackage{enumitem}
\usepackage{caption}
\usepackage{amsmath}
\usepackage{amsthm}
\usepackage{fontenc}

%% ---------------------------------------------------------------
%% Shared worksheet macros.
%%
%% These came from the Summer 2026 worksheet set, where each worksheet carried its
%% own copy. They have to live here instead: a xourse inlines every activity into a
%% single document, so duplicate \newcommand definitions across activities collide,
%% and the xourse gobbles \input so a per-topic macro file never loads.
%%
%% They use \providecommand, NOT \newcommand. ximera.cls loads this file automatically,
%% and every activity also carries an explicit \typeout{Start loading xmPreamble.tex}%

% Add here extra macro's that are loaded automatically by all documents of claas 'ximera' or 'xourse' in this repo

%%
%%  Example:
%%
% \newcommand{\R}{\mathbb{R}}
\usepackage{tikz}
\usetikzlibrary{positioning, arrows.meta}
\usepackage{multicol}
\usepackage{enumitem}
\usepackage{caption}
\usepackage{amsmath}
\usepackage{amsthm}
\usepackage{fontenc}

%% ---------------------------------------------------------------
%% Shared worksheet macros.
%%
%% These came from the Summer 2026 worksheet set, where each worksheet carried its
%% own copy. They have to live here instead: a xourse inlines every activity into a
%% single document, so duplicate \newcommand definitions across activities collide,
%% and the xourse gobbles \input so a per-topic macro file never loads.
%%
%% They use \providecommand, NOT \newcommand. ximera.cls loads this file automatically,
%% and every activity also carries an explicit \input{../xmPreamble}, so this file is
%% read twice. That was harmless while it held only \usepackage lines (LaTeX skips a
%% package it has already loaded) but a second \newcommand is a hard error.
%%
%%   \blank[width]        fill-in-the-blank rule (default 2.5cm)
%%   \showwork            "show and explain your work" reminder
%%   \showworkline        ... plus which schedule line was used
%%   \showworkyear        ... plus which year's schedule was used
%%   \schedhead{...}      centred bold caption above a tiered-rate table
%%   \samesquares[scale]{left}{right}   two equal-area squares, labelled
%%   \samecubes[scale]{left}{right}     two equal-volume cubes, labelled
%% ---------------------------------------------------------------

\providecommand{\blank}[1][2.5cm]{\underline{\hspace{#1}}}
\providecommand{\showwork}{\textbf{REMEMBER TO SHOW AND EXPLAIN YOUR WORK.}}
\providecommand{\showworkline}{\textbf{REMEMBER TO SHOW AND EXPLAIN YOUR WORK.
  Explanation should include how and why you know which line of the
  schedule was used to calculate this amount.}}
\providecommand{\showworkyear}{\begin{sloppypar}\textbf{REMEMBER TO SHOW AND
  EXPLAIN YOUR WORK. Explanation should include how and why you know which
  line of the year's tax schedule was used to calculate this tax
  amount.}\end{sloppypar}}
\providecommand{\schedhead}[1]{\begin{center}\textbf{#1}\end{center}}

\providecommand{\samesquares}[3][1]{%
\begin{center}
{\footnotesize These squares are the \textbf{same size}.}

\vspace{1.5mm}
\begin{tikzpicture}[scale=#1,every node/.style={scale=#1}]
  \draw (0,0) rectangle (2.4,2.4);
  \node[anchor=north west] at (0.15,2.25) {Area:};
  \node[anchor=east]  at (-0.15,1.2) {#2};
  \node[anchor=north] at (1.2,-0.15) {#2};
  \begin{scope}[xshift=4.6cm]
    \draw (0,0) rectangle (2.4,2.4);
    \node[anchor=north west] at (0.15,2.25) {Area:};
    \node[anchor=east]  at (-0.15,1.2) {#3};
    \node[anchor=north] at (1.2,-0.15) {#3};
  \end{scope}
\end{tikzpicture}
\end{center}}

\providecommand{\samecubes}[3][1]{%
\begin{center}
{\footnotesize These cubes are the \textbf{same size}.}

\vspace{1.5mm}
\begin{tikzpicture}[scale=#1,every node/.style={scale=#1}]
  \foreach \dx in {0,5.2} {
    \begin{scope}[xshift=\dx cm]
      \draw (0,0) rectangle (2.0,2.0);
      \draw (0,2.0) -- (0.65,2.6) -- (2.65,2.6) -- (2.0,2.0);
      \fill[black!12] (2.0,0) -- (2.65,0.6) -- (2.65,2.6) -- (2.0,2.0) -- cycle;
      \draw (2.0,0) -- (2.65,0.6) -- (2.65,2.6) -- (2.0,2.0);
      \node[anchor=north west] at (0.1,1.9) {Volume:};
    \end{scope}
  }
  \node[anchor=east]  at (-0.15,1.0) {#2};
  \node[anchor=north] at (1.0,-0.15) {#2};
  \node[anchor=west]  at (2.25,0.4) {#2};
  \node[anchor=east]  at (5.05,1.0) {#3};
  \node[anchor=north] at (6.2,-0.15) {#3};
  \node[anchor=west]  at (7.45,0.4) {#3};
\end{tikzpicture}
\end{center}}

%% NOTE: do NOT add \usepackage to individual activity .tex files.
%% When a xourse inlines an activity it gobbles \documentclass and \input, but a bare
%% \usepackage still executes -- after \begin{document} -- and errors out.
%% All shared packages and macros belong here.
%%
%% ximera.cls already loads: amsmath, amssymb, amsthm, cancel, enumitem, graphicx,
%% hyperref, listings, pgfplots, tikz, url, xcolor. No need to re-load those.
, so this file is
%% read twice. That was harmless while it held only \usepackage lines (LaTeX skips a
%% package it has already loaded) but a second \newcommand is a hard error.
%%
%%   \blank[width]        fill-in-the-blank rule (default 2.5cm)
%%   \showwork            "show and explain your work" reminder
%%   \showworkline        ... plus which schedule line was used
%%   \showworkyear        ... plus which year's schedule was used
%%   \schedhead{...}      centred bold caption above a tiered-rate table
%%   \samesquares[scale]{left}{right}   two equal-area squares, labelled
%%   \samecubes[scale]{left}{right}     two equal-volume cubes, labelled
%% ---------------------------------------------------------------

\providecommand{\blank}[1][2.5cm]{\underline{\hspace{#1}}}
\providecommand{\showwork}{\textbf{REMEMBER TO SHOW AND EXPLAIN YOUR WORK.}}
\providecommand{\showworkline}{\textbf{REMEMBER TO SHOW AND EXPLAIN YOUR WORK.
  Explanation should include how and why you know which line of the
  schedule was used to calculate this amount.}}
\providecommand{\showworkyear}{\begin{sloppypar}\textbf{REMEMBER TO SHOW AND
  EXPLAIN YOUR WORK. Explanation should include how and why you know which
  line of the year's tax schedule was used to calculate this tax
  amount.}\end{sloppypar}}
\providecommand{\schedhead}[1]{\begin{center}\textbf{#1}\end{center}}

\providecommand{\samesquares}[3][1]{%
\begin{center}
{\footnotesize These squares are the \textbf{same size}.}

\vspace{1.5mm}
\begin{tikzpicture}[scale=#1,every node/.style={scale=#1}]
  \draw (0,0) rectangle (2.4,2.4);
  \node[anchor=north west] at (0.15,2.25) {Area:};
  \node[anchor=east]  at (-0.15,1.2) {#2};
  \node[anchor=north] at (1.2,-0.15) {#2};
  \begin{scope}[xshift=4.6cm]
    \draw (0,0) rectangle (2.4,2.4);
    \node[anchor=north west] at (0.15,2.25) {Area:};
    \node[anchor=east]  at (-0.15,1.2) {#3};
    \node[anchor=north] at (1.2,-0.15) {#3};
  \end{scope}
\end{tikzpicture}
\end{center}}

\providecommand{\samecubes}[3][1]{%
\begin{center}
{\footnotesize These cubes are the \textbf{same size}.}

\vspace{1.5mm}
\begin{tikzpicture}[scale=#1,every node/.style={scale=#1}]
  \foreach \dx in {0,5.2} {
    \begin{scope}[xshift=\dx cm]
      \draw (0,0) rectangle (2.0,2.0);
      \draw (0,2.0) -- (0.65,2.6) -- (2.65,2.6) -- (2.0,2.0);
      \fill[black!12] (2.0,0) -- (2.65,0.6) -- (2.65,2.6) -- (2.0,2.0) -- cycle;
      \draw (2.0,0) -- (2.65,0.6) -- (2.65,2.6) -- (2.0,2.0);
      \node[anchor=north west] at (0.1,1.9) {Volume:};
    \end{scope}
  }
  \node[anchor=east]  at (-0.15,1.0) {#2};
  \node[anchor=north] at (1.0,-0.15) {#2};
  \node[anchor=west]  at (2.25,0.4) {#2};
  \node[anchor=east]  at (5.05,1.0) {#3};
  \node[anchor=north] at (6.2,-0.15) {#3};
  \node[anchor=west]  at (7.45,0.4) {#3};
\end{tikzpicture}
\end{center}}

%% NOTE: do NOT add \usepackage to individual activity .tex files.
%% When a xourse inlines an activity it gobbles \documentclass and \input, but a bare
%% \usepackage still executes -- after \begin{document} -- and errors out.
%% All shared packages and macros belong here.
%%
%% ximera.cls already loads: amsmath, amssymb, amsthm, cancel, enumitem, graphicx,
%% hyperref, listings, pgfplots, tikz, url, xcolor. No need to re-load those.
, so this file is
%% read twice. That was harmless while it held only \usepackage lines (LaTeX skips a
%% package it has already loaded) but a second \newcommand is a hard error.
%%
%%   \blank[width]        fill-in-the-blank rule (default 2.5cm)
%%   \showwork            "show and explain your work" reminder
%%   \showworkline        ... plus which schedule line was used
%%   \showworkyear        ... plus which year's schedule was used
%%   \schedhead{...}      centred bold caption above a tiered-rate table
%%   \samesquares[scale]{left}{right}   two equal-area squares, labelled
%%   \samecubes[scale]{left}{right}     two equal-volume cubes, labelled
%% ---------------------------------------------------------------

\providecommand{\blank}[1][2.5cm]{\underline{\hspace{#1}}}
\providecommand{\showwork}{\textbf{REMEMBER TO SHOW AND EXPLAIN YOUR WORK.}}
\providecommand{\showworkline}{\textbf{REMEMBER TO SHOW AND EXPLAIN YOUR WORK.
  Explanation should include how and why you know which line of the
  schedule was used to calculate this amount.}}
\providecommand{\showworkyear}{\begin{sloppypar}\textbf{REMEMBER TO SHOW AND
  EXPLAIN YOUR WORK. Explanation should include how and why you know which
  line of the year's tax schedule was used to calculate this tax
  amount.}\end{sloppypar}}
\providecommand{\schedhead}[1]{\begin{center}\textbf{#1}\end{center}}

\providecommand{\samesquares}[3][1]{%
\begin{center}
{\footnotesize These squares are the \textbf{same size}.}

\vspace{1.5mm}
\begin{tikzpicture}[scale=#1,every node/.style={scale=#1}]
  \draw (0,0) rectangle (2.4,2.4);
  \node[anchor=north west] at (0.15,2.25) {Area:};
  \node[anchor=east]  at (-0.15,1.2) {#2};
  \node[anchor=north] at (1.2,-0.15) {#2};
  \begin{scope}[xshift=4.6cm]
    \draw (0,0) rectangle (2.4,2.4);
    \node[anchor=north west] at (0.15,2.25) {Area:};
    \node[anchor=east]  at (-0.15,1.2) {#3};
    \node[anchor=north] at (1.2,-0.15) {#3};
  \end{scope}
\end{tikzpicture}
\end{center}}

\providecommand{\samecubes}[3][1]{%
\begin{center}
{\footnotesize These cubes are the \textbf{same size}.}

\vspace{1.5mm}
\begin{tikzpicture}[scale=#1,every node/.style={scale=#1}]
  \foreach \dx in {0,5.2} {
    \begin{scope}[xshift=\dx cm]
      \draw (0,0) rectangle (2.0,2.0);
      \draw (0,2.0) -- (0.65,2.6) -- (2.65,2.6) -- (2.0,2.0);
      \fill[black!12] (2.0,0) -- (2.65,0.6) -- (2.65,2.6) -- (2.0,2.0) -- cycle;
      \draw (2.0,0) -- (2.65,0.6) -- (2.65,2.6) -- (2.0,2.0);
      \node[anchor=north west] at (0.1,1.9) {Volume:};
    \end{scope}
  }
  \node[anchor=east]  at (-0.15,1.0) {#2};
  \node[anchor=north] at (1.0,-0.15) {#2};
  \node[anchor=west]  at (2.25,0.4) {#2};
  \node[anchor=east]  at (5.05,1.0) {#3};
  \node[anchor=north] at (6.2,-0.15) {#3};
  \node[anchor=west]  at (7.45,0.4) {#3};
\end{tikzpicture}
\end{center}}

%% NOTE: do NOT add \usepackage to individual activity .tex files.
%% When a xourse inlines an activity it gobbles \documentclass and \input, but a bare
%% \usepackage still executes -- after \begin{document} -- and errors out.
%% All shared packages and macros belong here.
%%
%% ximera.cls already loads: amsmath, amssymb, amsthm, cancel, enumitem, graphicx,
%% hyperref, listings, pgfplots, tikz, url, xcolor. No need to re-load those.

\title{In-Class: Percentages}
\begin{document}
\begin{abstract}
Percentages as relative comparisons: percent increase and decrease, flat taxes
expressed as percentages, and several situations where percentages are commonly
misinterpreted --- comparing percentages of different wholes, and chaining percent
changes.
\end{abstract}
\maketitle

\section*{Percentages}

A very common use of percentages is as a \emph{relative comparison} value. Using
percents this way can put value changes in perspective. For example, if the price of
something increases from \$0.50 to \$1, that is ``only'' an absolute increase of
\$0.50, which may not seem like much --- but that is a relative increase of $100\%$!
Similarly, if a construction project goes \$10,000 over budget, that may initially
seem like a lot, but if the original budget was \$1 million then they only went $1\%$
over budget.

Let's look at a few problems which use percentages to make relative comparisons.

\subsection*{Section 1: Percent change}

The recipe throughout is the same two steps:
\begin{enumerate}
\item Find the \emph{absolute} change: new value $-$ original value.
\item Express that absolute change as a percentage \emph{of the original value}.
\end{enumerate}

\begin{example}
Target increased the price of their Cheerios from \$2.99 to \$3.64. By what percent
did they increase the price?

Absolute price change: $\$3.64 - \$2.99 = \$0.65$.

What is the absolute change as a percentage of the original price? That is, \$0.65 is
what percent of the original price? So $\$0.65 = x \cdot \$2.99$, where $x$ is the
unknown percentage. Solving,
\[
x = \frac{\$0.65}{\$2.99} \approx 0.21739,
\]
and $0.21739 \cdot 100 \approx 21.74$, so the price increased by $21.74\%$.
\end{example}

\begin{problem}
Back in 2016, a company which sells EpiPens decided to boost the prices. In a very
short period of time, the price of a pair of EpiPens increased from about \$115 to
about \$610. What was the percent increase in the price?

The absolute change in price is $\$\answer{495}$, which is a
$\answer[tolerance=0.5]{430.4}$\% increase.

\begin{solution}
The absolute change is $\$610 - \$115 = \$495$. Now express it as a percent of the
original price: $\$495 = x \cdot \$115$, so
\[
x = \frac{\$495}{\$115} \approx 4.304,
\]
which is $430.4\%$. So the price increase from \$115 to \$610 is a $430.4\%$ increase.
\end{solution}
\end{problem}

\begin{problem}
The tax on gasoline has been $18.4$ cents per gallon since 1993, the last tax
increase. In 1993 a gallon of gas cost about \$1.16 (including the tax). In 2011 gas
was about \$3.40 per gallon (still including the tax).

\begin{prompt}
\textbf{(a)} What are the 1993 and 2011 costs per gallon \emph{not} including tax?

1993: $\$\answer{0.976}$ \qquad 2011: $\$\answer{3.216}$

\begin{hint}
$18.4$ cents $= \$0.184$. Subtract this tax off in each year.
\end{hint}
\end{prompt}

\begin{prompt}
\textbf{(b)} What is the $18.4$ cents as a percent tax in 1993?
$\answer[tolerance=0.05]{18.85}$\%

\begin{solution}
Express \$0.184 as a percentage of the untaxed 1993 cost:
$\$0.184 = x \cdot \$0.976$ gives $x = \frac{\$0.184}{\$0.976} \approx 0.18852$,
which is about $18.85\%$.
\end{solution}
\end{prompt}

\begin{prompt}
\textbf{(c)} What is the $18.4$ cents as a percent tax in 2011?
$\answer[tolerance=0.05]{5.72}$\%

\begin{solution}
Same argument, but with the untaxed 2011 cost:
$\$0.184 = x \cdot \$3.216$ gives $x = \frac{\$0.184}{\$3.216} \approx 0.0572$,
which is about $5.72\%$.

The flat tax has not changed at all, but as a share of the price it has shrunk
dramatically.
\end{solution}
\end{prompt}

\begin{prompt}
\textbf{(d)} How much would the 2011 gallon of gas cost if, instead of a flat rate of
$18.4$ cents per gallon, the tax were the same \emph{percent} tax you calculated for
1993?

$\$\answer[tolerance=0.02]{3.82}$ per gallon.

\begin{solution}
The untaxed 2011 cost is \$3.216 and the 1993 percent tax is $18.85\%$, so the tax
would be $0.1885 \cdot \$3.216 \approx \$0.606$. The cost with tax included would be
$\$3.216 + \$0.606 = \$3.822$, about \$3.82 per gallon.
\end{solution}
\end{prompt}

\begin{prompt}
\textbf{(e)} What was the percent increase in gas prices from 1993 to 2011?
$\answer[tolerance=0.5]{193.1}$\%

Should you include the tax to calculate this percent change? Why or why not?
\begin{freeResponse}
\end{freeResponse}

\begin{solution}
Looking at the overall price with tax included, the absolute change is
$\$3.40 - \$1.16 = \$2.24$. As a percentage of the original cost,
$\$2.24 = x \cdot \$1.16$, so $x = \frac{\$2.24}{\$1.16} \approx 1.931$, about
$193.1\%$. So there is a $193.1\%$ increase in gas prices from 1993 to 2011.
\end{solution}
\end{prompt}
\end{problem}

\subsection*{Section 2: Percentages that get misread}

Below are some situations involving percentages that are often incorrectly
interpreted. Let's work through these problems (hopefully correctly) and reflect on
what mistaken interpretations may come up.

\begin{problem}
A newspaper headline claims: ``Apple spends far less on R\&D than any of its rivals.''
Further into the story, the reporter writes that Apple spends a paltry $2\%$ of
revenues compared with $14\%$ at Google and Microsoft.

\begin{prompt}
\textbf{(a)} Give example numbers where $2\%$ of some amount is actually larger than
$14\%$ of another amount. (You answered this in your pre-class work --- see what the
others in your group used.)
\begin{freeResponse}
\end{freeResponse}
\end{prompt}

\begin{prompt}
\textbf{(b)} Why can we not say that the $2\%$ in this situation is definitely less
than the $14\%$?
\begin{freeResponse}
\end{freeResponse}
\end{prompt}

\begin{prompt}
\textbf{(c)} Let Apple's revenue be $A$ and Google's revenue be $G$.

\textbf{i.} Write how you would calculate $2\%$ of $A$.
\begin{freeResponse}
\end{freeResponse}

\textbf{ii.} Write how you would calculate $14\%$ of $G$.
\begin{freeResponse}
\end{freeResponse}

\textbf{iii.} How much larger can Apple's revenue be while $2\%$ of it is still more
than $14\%$ of Google's revenue? (Hint: set up an inequality using parts i and ii.)
\begin{freeResponse}
\end{freeResponse}

\textbf{iv.} Would $2\%$ of Apple's 2017 revenue be more than $14\%$ of Google's 2017
revenue? Look it up!
\begin{freeResponse}
\end{freeResponse}

\begin{solution}
From $0.02A > 0.14G$ we get $A > 7G$: Apple's revenue would have to be more than
seven times Google's for the comparison to flip.
\end{solution}
\end{prompt}
\end{problem}

\begin{problem}
A shop owner raises the price of a \$100 pair of shoes by $50\%$. After a few weeks,
because of falling sales, the owner reduces the price of the shoes by $50\%$.

\begin{prompt}
\textbf{(a)} What is the new price of the shoes, after both percent changes?
$\$\answer{75}$

\begin{solution}
Increase: $\$100 + 0.50 \cdot \$100 = \$150$.
Decrease: $\$150 - 0.50 \cdot \$150 = \$150 - \$75 = \$75$.
\end{solution}
\end{prompt}

\begin{prompt}
\textbf{(b)} Why are the shoes not back at the \$100 price? Don't say something like
``because of the math above'' --- try to explain why the math doesn't return the
original amount.
\begin{freeResponse}
\end{freeResponse}

\begin{solution}
The $50\%$ decrease is $50\%$ of the \emph{current} price, \$150, while the original
$50\%$ increase was $50\%$ of the \emph{original} price, \$100. Percentages are
relative --- it matters what whole you are taking the percentage of.
\end{solution}
\end{prompt}

\begin{prompt}
\textbf{(c)} How much would the shop owner have to decrease the price in order to
have the shoes back at \$100 after the $50\%$ increase?
$\answer[tolerance=0.05]{33.33}$\%

\begin{solution}
After the increase the price is \$150, and we need to come down by \$50. Express \$50
as a percentage of the current price: $\$50 = x \cdot \$150$, so
$x = \frac{50}{150} = 0.333\ldots$, about $33.33\%$.
\end{solution}
\end{prompt}
\end{problem}

\begin{problem}
The annual number of burglaries in a town rose by $40\%$ in 2012 and fell by $30\%$
in 2013.

\begin{prompt}
\textbf{(a)} What was the total percent change in burglaries over the two years?

First try the problem with a specific starting number. In your group, pick a 3-digit
number that does NOT end in $0$ --- say $355$ burglaries in 2011 --- and find the
percent change. Then think about how to do the problem without knowing a specific
starting number.

Starting from $355$: there were $\answer{497}$ burglaries in 2012 and about
$\answer[tolerance=1]{348}$ in 2013, an overall
$\answer[tolerance=0.05]{1.97}$\% decrease.

\begin{solution}
2012: $355 + 0.40 \cdot 355 = 355 + 142 = 497$ burglaries.

2013: $497 - 0.30 \cdot 497 \approx 348$ burglaries.

Compare original to final. Absolute change: $355 - 348 = 7$. Now $7$ is what percent
of the original? $7 = x \cdot 355$, so $x = \frac{7}{355} \approx 0.0197$, about
$1.97\%$. So there is a $1.97\%$ \emph{decrease} in burglaries from 2011 to 2013.

Without a specific starting number: a $40\%$ rise then a $30\%$ fall multiplies the
original by $1.40 \cdot 0.70 = 0.98$ --- a $2\%$ decrease regardless of where you
start.
\end{solution}
\end{prompt}

\begin{prompt}
\textbf{(b)} It might be tempting to say the change over the two years was a $10\%$
increase. Why might someone think that --- how do you get $10\%$ from the numbers in
the problem?
\begin{freeResponse}
\end{freeResponse}

\begin{solution}
Someone might simply compute $40\% - 30\% = 10\%$.
\end{solution}
\end{prompt}

\begin{prompt}
\textbf{(c)} Why is $10\%$ incorrect?
\begin{freeResponse}
\end{freeResponse}

\begin{solution}
The $40\%$ and the $30\%$ are applied to different quantities --- they are
proportions of different wholes. The $40\%$ is $40\%$ of the 2011 burglary count,
while the $30\%$ is $30\%$ of the 2012 count. Since the percentages are relative to
different values, we cannot subtract them or compare them directly.
\end{solution}
\end{prompt}
\end{problem}

\end{document}
